\subsection{Arithmetic nonvanishing from the integral class at two}

\begin{statement}{B.37}{The actual outer extension and the arithmetic Ross edge are rationally nonzero}{Proved}
Use the actual $C_0$ tangent and integral class $\eta$ of
B.31--B.36. Put
\[
 Q=\operatorname{Lie}_3(V)(-1)/\operatorname{ad}_w(V),
 \qquad \overline\eta\in H^1(\Gamma,Q),
 \qquad
 \Gamma=G_{\mathbb Q(\zeta_3),\{3\}}(3)\rtimes\langle c\rangle.
                                                               \tag{B.37.1}
\]
Then $Q$ is a free rank-$320$ integral lattice and
\[
       [\overline\eta]\ne0
          \quad\hbox{in }H^1(\Gamma,Q\otimes\mathbb Q_3).
                                                               \tag{B.37.2}
\]
More precisely, none of the $27$ local candidates of B.36,
after projection to $Q$, is the restriction of a rationally
zero global class over $K=\mathbb Q(\zeta_3)$.
An explicit functional with five nonzero Hall coefficients
modulo nine proves this exclusion.

Let $D$ and $\Phi:D\xrightarrow{\sim}H(2)$ be the
specified two-dimensional invariant component and
normalization of B.20. The rational class of $\overline\eta$
is the image of the $D(-1)$ component of $\eta$.
Consequently
\[
 \Phi_*[\eta]_{\rm ell}
    =-\iota_{w*}\operatorname{edge}_H
                 \bigl(\kappa(1-x)\smile\kappa(1-y)\bigr)
       \ne0\quad\hbox{in }H^1(\Gamma,H(1)).                  \tag{B.37.3}
\]
The last arithmetic cohomology space has dimension one
over $\mathbb Q_3$, proved below directly for this
bad-prime-three source. Thus the specified arithmetic
Ross edge is a generator of that rational vector space.

This proves the arithmetic nonvanishing clause of prompt
B.4. The integral refinement is a separate question from
this rational nonvanishing argument and from the order
$m_2=3$. B.38 computes its exact homogeneous ambiguity,
and B.40 selects the actual phase by evaluating the
geometric extension and the intrinsic finite detector.
\end{statement}
\begin{proof}
\emph{A necessary local condition for rationally zero classes.}
Let $T$ be any finite free $\mathbb Z_3$-module with a
continuous arithmetic action. Suppose that two arithmetic
elements $t,F$ commute on $T$, and that $F-I$ is
invertible on $T\otimes\mathbb Q_3$.
If an integral cocycle $z$ is rationally a coboundary,
write $z(g)=(g-I)y$ for one
$y\in T\otimes\mathbb Q_3$. Its local unramified class
at $F$ is represented by $(F-I)y$ in $T/(F-I)T$.
Since $z(t)=(t-I)y$ is integral,
\[
 (t-I)z(F)=(F-I)(t-I)y\in(F-I)T.                            \tag{B.37.4}
\]
Hence this local class is fixed by $t$.
This argument uses actual coefficient matrices only;
it does not require a word expressing $F$ in global
generators. If a set of such $t$ generates the arithmetic
action, the converse holds as well: for a fixed local
class represented by $a$, put $y=(F-I)^{-1}a$.
The fixed-class conditions say $(t-I)y\in T$ on the
generators, hence on their closed generated group.
Thus the fixed subgroup is exactly the local image of
the rationally zero global classes. We need only the
necessary implication (B.37.4).

\emph{The actual global Frobenii in the local coefficient basis.}
All nine cubic cusps are rational over $K$. At seven
their reductions give all nine points of the compact
curve: every nonzero cube in $\mathbb F_7$ is $1$ or
$-1$, and three such cubes cannot sum to zero.
Thus its compact Frobenius polynomial is $X^2+X+7$.
Let $S$ be the compatible CM matrix of B.35. Its two
roots in the CM endomorphism ring are $I+3S$ and
$I+3S^2$. The two actual Frobenii above seven have these
compact matrices, in either order.
Their full degree-one matrices, and that of Frobenius
four at two, are therefore
\[
 \begin{aligned}
 T_7&=\begin{pmatrix}I+3S&0\\U(S-2I)&7I_8\end{pmatrix},
 &T'_7&=\begin{pmatrix}I+3S^2&0\\U(S^2-2I)&7I_8\end{pmatrix},\\
 F_4&=\begin{pmatrix}-2I_2&0\\-2U&4I_8\end{pmatrix}.
 \end{aligned}                                                \tag{B.37.5}
\]
The cross blocks follow from the actual gluing formula
(B.24.3), not from a chosen direct-sum replacement.
The three matrices commute. On $A(-1)$ their actions
are $7^{-1}\operatorname{Lie}_3(T_7)$,
$7^{-1}\operatorname{Lie}_3(T'_7)$, and
$4^{-1}\operatorname{Lie}_3(F_4)$.
They preserve $\operatorname{ad}_w(V)$ and induce
commuting actions on $Q$. We have passed to $K$, so
no unidentified change of basis for global complex
conjugation is combined with the local quaternion basis.

The primitivity assertion of B.35 also has a literal
integral quotient basis. Order $V$ as
$(e_0,\ldots,e_9)=(f_1,f_2,A_0,A_1,A_2,B_0,B_1,B_2,C_1,C_2)$.
The ten columns
\[
             (u,v,e_0,e_1,e_3,e_5,e_6,e_7,e_8,e_9)
                                                               \tag{B.37.6}
\]
have determinant $-1$.
Write $e'_0,\ldots,e'_9$ for this adapted basis.
In its associated Hall basis the ten vectors
$\operatorname{ad}_w(e'_j)$ are exactly the distinct
columns indexed by $(1,0,j)$, $0\le j\le9$.
Deleting those columns gives $Q$, with no saturation
ambiguity. The induced Frobenius-four-minus-one operator
has Smith valuations
\[
                     0^{122},\ 1^{70},\ 2^{26},\ 3^{82},\ 4^{20}.
                                                               \tag{B.37.7}
\]
Every one of the $27$ projected classes from B.36 still
has order $27$. Its difference under either seven
Frobenius has exact order $9$ in this cokernel.
The following smaller certificate proves the nonzero
differences without relying only on these Smith counts.

\emph{A five-entry certificate excluding all global torsion.}
Continue to use the original ordered basis $e_0,\ldots,e_9$
and the Hall convention $[[e_i,e_j],e_k]$, with $i>j$
and $k\ge j$. Define $\ell:A\to\mathbb Z/9$ by
\[
 \begin{array}{c|rrrrr}
 (i,j,k)&(2,0,4)&(4,0,2)&(4,0,4)&(4,0,8)&(8,0,4)\\ \hline
 \ell([[e_i,e_j],e_k])&-1&-2&3&1&-4
 \end{array}                                                  \tag{B.37.8}
\]
and set it equal to zero on every other Hall vector.
Put $D_4=4^{-1}\operatorname{Lie}_3(F_4)-I$.
Direct integral matrix multiplication modulo nine gives
\[
             \ell D_4=0,\qquad \ell\operatorname{ad}_w=0.
                                                               \tag{B.37.9}
\]
Thus $\ell$ is a well-defined functional on
$A/(D_4A+\operatorname{ad}_wV)=Q/(F_4-I)Q$.
Let $B$ be the restriction matrix (B.36.16), and write
$x_1,\ldots,x_{13}$ for the Smith coordinates (B.36.15).
The same multiplication, now using (B.37.5), gives
\[
 \ell\bigl((7^{-1}\operatorname{Lie}_3(T_7)-I)Bx\bigr)
       =3x_7+6(x_8+x_9+x_{10}+x_{11})\pmod9.                 \tag{B.37.10}
\]
For $T'_7$ the right side is its negative.
All these equations are finite identities in explicitly
given matrices; no truncated Kummer logarithm is treated
as an exact three-adic number.

Evaluate the right side on the complete parametrization
$b+s h_3+t h_9$ of B.36.18--B.36.19. On $b$ it is
$3\cdot5+6(9+6+9+24)=303\equiv6\pmod9$.
On both $h_3$ and $h_9$ it is zero modulo nine.
Therefore it equals $6$ on all $27$ candidates, and
the conjugate Frobenius gives $3$ on all of them.
None is fixed by either actual seven Frobenius.
Equation (B.37.4) proves that the actual integral
$\overline\eta$ is not rationally zero over $K$,
and hence cannot be rationally zero over $\mathbb Q$.
This proves (B.37.2).

There are two separately executed coefficient calculations.
The program \path{c_char2_inner_quotient.py} constructs
the unimodular basis (B.37.6), literally deletes the ten
inner Hall columns, and computes the rank-$320$ matrices.
The program \path{b_char2_outer_torsion.py} instead uses
the rectangular presentation $(D_4,\operatorname{ad}_w)$
on the original $330$ rows; its Smith valuations are
those of (B.37.7) with ten additional unit factors.
Both verify all $27$ differences and the five-entry
functional (B.37.8)--(B.37.10).
The original matrix precision is $3^{10}$, whereas this
certificate is modulo nine and every local Smith pivot
is at most $3^4$. Thus no uncomputed higher digit enters
the exclusion.

\emph{Why the detected rational class is the Ross component.}
Over $\mathbb Q_3$ write $V=H\oplus W$ as in B.20.
This is the splitting of the generalized-Jacobian
extension whose cusp divisor classes are three-torsion.
The torsion cusp divisors give this Galois-stable
rational splitting. Its two summands have distinct
weights at a good Frobenius.
Each automorphism of the marked curve over
$K$ commutes with that Frobenius on $V$, so preserves
this splitting as well.

Use a rational class-three Malcev Lie algebra and choose
the pure degree-two lift $\omega=[u,v]$ with $u,v\in H$.
This is a permitted rational change of the lift defining
$\eta$; its difference from an integral Hall lift adds a
rational coboundary. We use the pure Lie bracket, not
the logarithm of a group commutator with an unremoved
cubic term. Write the action on a degree-one generator
as $g(a)=A_ga+B_ga+C_ga$ by degrees one, two, and three.
Modulo degree four,
\[
 g[u,v]-\chi_3(g)[u,v]
             =[B_gu,A_gv]+[A_gu,B_gv].                       \tag{B.37.11}
\]
The degree-three generator terms contribute only in
degree four. Since $A_gu,A_gv\in H$ and
$B_g(H)\subset\Lambda^2(H\oplus W)$, the cocycle
therefore lies in
\[
       \operatorname{Lie}_3(H)
         +(H^2W)+J(H\otimes\Lambda^2W).                      \tag{B.37.12}
\]
Here $H^2W$ denotes the full multidegree with two
compact and one boundary letter. There is no pure
$W^3$ contribution and no other $HW^2$ contribution.

Let $\mathcal A=E[3]\rtimes\mu_6$ be the group of $54$
global Fermat automorphisms used in B.20.
All its elements are defined over $K$, fix $w$, and
satisfy the automorphism/path identity (B.36.2).
Thus the class of $\overline\eta$ over $K$ is
$\mathcal A$-invariant. Individual automorphisms need
not commute with $G_{\mathbb Q}$; we make this assertion
over $K$. Their rational average
$e_{\mathcal A}=54^{-1}\sum_{a\in\mathcal A}a$ does
commute with $G_{\mathbb Q}$, because Galois permutes
the $54$ automorphisms. Restriction from $\mathbb Q$ to
$K$ is injective in rational cohomology, as its
corestriction composite is multiplication by two.
Consequently
\[
                  e_{\mathcal A}[\overline\eta]
                         =[\overline\eta]
              \quad\hbox{over }\mathbb Q.                   \tag{B.37.13}
\]

The translation group $E[3]$ fixes $H$ and acts regularly
on the nine cusps. Its invariants in the reduced
permutation representation $W$ are zero. It therefore
kills the $H^2W$ term of (B.37.12) on averaging.
The origin-fixing involution $[-1]$ acts as $-1$ on
$\operatorname{Lie}_3(H)$, so kills that term.
The remaining invariant space is exactly
$D=(J(H\otimes\Lambda^2W))^{\mathcal A}$ of B.20.
Moreover $H^{\mu_6}=0$ and $W^{E[3]}=0$, whence
$V^{\mathcal A}=0$. Since $w$ is fixed,
\[
            e_{\mathcal A}\operatorname{ad}_w(V)
                        =\operatorname{ad}_w(V^{\mathcal A})=0.
                                                               \tag{B.37.14}
\]
Taking invariants is exact over $\mathbb Q_3$.
Thus $A(-1)^{\mathcal A}\to Q^{\mathcal A}$ is an
isomorphism, and the nonzero class in (B.37.13) is
precisely the image of the $D(-1)$ component of $\eta$.
Equations (B.37.11)--(B.37.14) exclude an alternative
Tate component from the detected rational class.
The signed comparison and the exact isomorphism $\Phi$
of B.20 now prove the nonvanishing in (B.37.3).

\emph{The arithmetic space really is one-dimensional at three.}
For completeness this calculation requires no good
reduction assumption at the coefficient prime.
Put $P=G_{K,\{3\}}(3)$ and
$X=\operatorname{Spec}\mathbb Z[\zeta_3,1/3]$.
The Eisenstein integer ring is principal, so
$\operatorname{Pic}(X)=0$.
The Brauer group of $X$ is zero: it injects into the
Brauer group of $K$, every local invariant except the
one prime above three vanishes, and reciprocity forces
that final invariant to vanish as well.
The degree-two injection for the universal pro-three
cover, proved explicitly in B.29, and the Kummer sequence
therefore give
\[
 H^2(P,\mathbb F_3)\lhook\joinrel\longrightarrow
 H^2_{\rm et}(X,\mu_3)=0.                                   \tag{B.37.15}
\]
Degree-one Kummer theory gives
$H^1(P,\mathbb F_3)=\langle\zeta_3,1-\zeta_3\rangle$,
of dimension two. The two Frobenii above seven have
the invertible evaluation matrix
$\left(\begin{smallmatrix}2&0\\2&2\end{smallmatrix}\right)$
from (B.29.5), so they topologically generate $P$.

Here is also the elementary freeness argument needed
for this chosen generating pair. Map the free pro-three
group on $\tau,\tau'$ onto $P$ and call its kernel $R$.
The isomorphism in $H^1(-,\mathbb F_3)$ and (B.37.15),
through the five-term sequence, give
$H^1(R,\mathbb F_3)^P=0$.
If $R$ were nontrivial, it would have nontrivial image
in a finite three-group quotient. The Frattini quotient
of that image has nonzero coinvariants under the finite
three-group, because its augmentation ideal is nilpotent.
A nonzero linear functional on those coinvariants would
give a nonzero invariant character of $R$, a contradiction.
Hence $P$ is free pro-three on this actual pair.
Choose the prolongations so that $c\tau c=\tau'$.

For the lattice $M=H_{\mathbb Z_3}(1)$ in its original
compact basis, the matrices at $\tau,c$ are
\[
 T_M=7\begin{pmatrix}1&-3\\3&-2\end{pmatrix},\qquad
 C_M=-\begin{pmatrix}1&-1\\0&-1\end{pmatrix}.                  \tag{B.37.16}
\]
Averaging by $c$ normalizes a cocycle to have value
zero at $c$. It is then determined by an arbitrary
value $a\in M$ at $\tau$, with value $C_Ma$ at $\tau'$.
Every such assignment is continuous: the residual
arithmetic action is a three-group, and the free
pro-three universal property supplies the section
into $M\rtimes P$ with the two assigned generator
values. The normalized coboundaries are precisely
$(T_M-I)M^{c=1}$.
The latter invariant lattice is $\mathbb Z_3(1,2)$, and
\[
 (T_M-I)\binom12=-9\binom41,\qquad
 H^1(\Gamma,M)=
       \mathbb Z_3^2/9\mathbb Z_3(4,1)
          \simeq\mathbb Z_3\oplus\mathbb Z/9.                \tag{B.37.17}
\]
The vector $(4,1)$ is primitive. Rationalization proves
the asserted dimension one, including at the bad
coefficient prime three. Formula (B.37.17) is also
checked directly by \path{b_ross_hodge.py}.
The nonzero class in (B.37.3) is consequently a rational
generator with the signed scalar of B.20.

Neither this dimension calculation nor the exclusion
modulo nine identifies the integral coordinate of that
generator in (B.37.17). A complete integral Frobenius
value and the geometrically admissible torsion data
remain to be determined. In particular the real value
$-3/2$ of B.19 continues to be the stated Hodge coordinate
of the Ross symbol; it is not being substituted for
an integral three-adic value of $\eta$.
\end{proof}
